% part: intuitionistic-logic
% chapter: soundness-completeness
% section: decidability

\documentclass[../../../include/open-logic-section]{subfiles}

\begin{document}

\olfileid{int}{sc}{dec}

\olsection{నిర్ణేయత}
$\Gamma \Proves/ !A$ అయిన ప్రతి సందర్భంలోనూ,
సంపూర్ణతా సిద్ధాంత నిరూపణ
$\Gamma \Entails/ !A$కు సాక్ష్యంగా
అనంత లోకాలు గల నమూనాను ఇస్తుందని
గమనించండి. కింది ప్రతిపాదన ప్రకారం,
$\Entails !A$ను నిరూపించడానికి అన్ని పరిమిత
నమూనాల్లో (అంటే పరిమిత లోకాల సమితి గల
నమూనాల్లో) $\mSat{M}{!A}$ అని
నిరూపిస్తే చాలు.

\begin{thm}\ollabel{thm:decidability}
  $\Entails/ !A$ అయితే
  $\mSat/{M'}{!A}$ అయ్యే పరిమిత నమూనా ఉంది.
\end{thm}

\begin{proof}
  $\mModel{M}=\tuple{W, R, V}$లో
  $\mSat/{M}{!A}$ అని, $P$ అనేది~$!A$లో
  కనిపించే \tetoken{ప్రతిజ్ఞావాక్య
  చరాల}{!!{propositional variable}s}
  సమితి అని, $S$ అనేది ఆ లక్ష్య సూత్రం యొక్క
  అన్ని ఉపసూత్రాల సమితి అని అనుకుందాం.
  $\mModel{M'}=\tuple{W', R', V'}$ను ఇలా
  నిర్వచిద్దాం: $W' = \Setabs{[w]}{w \in W}$,
  ఇక్కడ $[w] = \Setabs{!B \in S}{\mSat{M}{!B}[w]}$;
  $R'$ ఉపసమితి సంబంధం; మరియు
  $V'(p)=\Setabs{[w]}{p \in [w]}$.
  $W'$ పరిమిత సమితి అని,
  $\mModel{M'}$ \tetoken{ఒక సంబంధ
  నమూనా}{!!a{relational model}} అని
  స్పష్టమే.

  $!A$పై ఆగమన పద్ధతిలో,~$P$ నుంచి
  \tetoken{ప్రతిజ్ఞావాక్య చరాలు}{!!{propositional variable}s}
  మాత్రమే గల $S$లోని అన్ని
  \tetoken{సూత్రాల}{!!{formula}s}
  $!A$కు
  \[
    \mSat{M}{!A}[w] \text{ అయితే, అప్పుడే } \mSat{M'}{!A}[{[w]}]
  \]
  అని చూపవచ్చు. దీన్ని పాఠకునికి
  అభ్యాసంగా వదులుతున్నాం.
  \sourcecorrection{OLTEINTDEC-001}{మూల నిర్మాణం లోకాలను $P$లోని అణు చరాల సత్యంతో మాత్రమే ఏకీకరించి, వాటి ఆధారంగా అన్ని సూత్రాల సత్యం నిలుస్తుందని పేర్కొంది; సోపాధికానికి అది తప్పు. పరిమిత ఉపసూత్ర సమితి $S$లో సత్యం ద్వారా లోకాలను వేరు చేసి, ఆ సమితిలోని సూత్రాలకు మాత్రమే ఆగమన వాదనను వర్తింపజేశాం.}
\end{proof}

\begin{prob}
~$P$ నుంచి ప్రతిజ్ఞావాక్య చరాలు మాత్రమే
గల $S$లోని అన్ని \tetoken{సూత్రాల}{!!{formula}s}~$!A$కు
$\mSat{M,w}{!A}$ అయితే, అప్పుడే
$\mSat{M',[w]}{!A}$ అని చూపి
\olref[int][sc][dec]{thm:decidability}
నిరూపణను పూర్తి చేయండి.
\end{prob}

\olref{thm:decidability} ప్రకారం
~$\Entails !A$ నిజమో కాదో నిర్ణయించే
కలనం ఉంది.

\end{document}
